We hypothesise that information organisation based on faceted classification is useful to improve enterprise information retrieval systems. The existence of similar facets in documents from different companies and the known adaptability of facet organisation strengthen this hypothesis. We refer this work to the automated classification and indexing on large amounts of text files. This work is descriptive, applied, and experimental. It aimed to expose the main characteristics of the enterprise information, proposing a tentative generalisation to the enterprise domain and presenting some facets we can use to organise it and to support better information retrieval. It applied facet analysis to two enterprise collections and evaluated the resulting faceted classification. Terms were selected from documents and queries. We found twelve common categories and the distribution of document subjects across the categories presents strong positive correlation by the Spearman’s rank correlation. Then, we obtained ten user queries and we adopted them to validate the found categories. We also used the Enterprise track of Text Retrieval Conference and its previous results as a Cranfield-like evaluation. The automated prototype used spatial, temporal, document and social characteristics. Thus, our empirical evaluation improved the information retrieval with no external dependency like Wikipedia or metasearch engines. The facet analysis was useful for comparing the companies with no desire to expose their information. The method can guide and stimulate future work and other companies can become more willing to take part in a research study.